\documentclass[11pt]{article}

\usepackage[margin=0.9in]{geometry}
\usepackage{amsmath,amssymb}
\usepackage{booktabs}
\usepackage{enumitem}
\usepackage{listings}
\usepackage{xcolor}
\usepackage{hyperref}

\definecolor{backcolour}{rgb}{0.95,0.95,0.92}
\definecolor{codeblue}{rgb}{0.05,0.15,0.55}
\definecolor{codegreen}{rgb}{0.0,0.45,0.0}

\lstset{
    backgroundcolor=\color{backcolour},
    basicstyle=\ttfamily\small,
    keywordstyle=\color{codeblue}\bfseries,
    commentstyle=\color{codegreen},
    stringstyle=\color{purple},
    breaklines=true,
    frame=single,
    showstringspaces=false,
    tabsize=4
}

\title{CPSC 250 \\ Recursion, Binary Search, and Algorithm Efficiency}
\author{Edward Brash}
\date{}

\begin{document}

\maketitle

\section{Recursion}

A recursive function is a function that calls itself.

Recursion can be useful when problems naturally break into smaller versions of themselves.

\section{A Bad Recursive Example}

Consider:

\begin{lstlisting}[language=Python]
def count_down(count):

    count_down(count - 1)
\end{lstlisting}

What happens if we call:

\begin{lstlisting}[language=Python]
count_down(3)
\end{lstlisting}

The function repeatedly calls itself forever:

\[
3 \rightarrow 2 \rightarrow 1 \rightarrow 0 \rightarrow -1 \rightarrow \cdots
\]

Eventually Python crashes with a recursion error.

\section{The Importance of an End Condition}

Recursive functions require an end condition.

\begin{lstlisting}[language=Python]
def count_down(count):

    if count == 0:
        print("Go!")

    else:
        print(count)
        count_down(count - 1)
\end{lstlisting}

Now:

\[
count\_down(3)
\]

produces:

\[
3, 2, 1, Go!
\]

\section{Recursion vs Iteration}

Although recursion can be elegant, it is often:

\begin{itemize}
    \item harder to understand,
    \item harder to debug,
    \item not more efficient.
\end{itemize}

One major exception is binary-search-style algorithms.

\section{Binary Search Motivation}

Suppose we want to find the root of a function:

\[
f(x) = 0
\]

Assume we know the root lies somewhere between:

\[
x_{low}
\]

and:

\[
x_{high}
\]

\section{Root Condition}

A root exists between two points if:

\[
f(x_{low}) \cdot f(x_{high}) < 0
\]

This means the function changes sign.

\section{The Binary Search Idea}

Suppose the root lies between:

\[
a
\]

and:

\[
b
\]

We compute the midpoint:

\[
c = \frac{a+b}{2}
\]

Then determine which half contains the root.

\section{Binary Search Algorithm}

\begin{enumerate}
    \item Compute midpoint:
    \[
    c = \frac{a+b}{2}
    \]

    \item Evaluate:
    \[
    f(a)f(c)
    \]

    \item If:
    \[
    f(a)f(c) < 0
    \]
    then the root lies between \(a\) and \(c\).

    Otherwise it lies between \(c\) and \(b\).

    \item Repeat.
\end{enumerate}

\section{Stopping Condition}

The algorithm ends when the interval becomes sufficiently small:

\[
|b-a| < \epsilon
\]

where:

\[
\epsilon
\]

is a tiny number.

\section{Example Function}

\begin{lstlisting}[language=Python]
def f(x):
    return -4*x + 3
\end{lstlisting}

This function has a root at:

\[
x = 0.75
\]

\section{Starting Interval}

\begin{lstlisting}[language=Python]
a = 0.0
b = 5.0
\end{lstlisting}

Check:

\begin{lstlisting}[language=Python]
f(a) * f(b)
\end{lstlisting}

If the product is negative, a root exists on the interval.

\section{Recursive Binary Search Root Finder}

\begin{lstlisting}[language=Python]
def find_root(a, b):

    epsilon = 1.0E-5

    if abs(b - a) < epsilon:
        print("Found root at x =", a)

    else:

        c = (a + b) / 2.0

        if f(a) * f(c) < 0:
            find_root(a, c)

        else:
            find_root(c, b)
\end{lstlisting}

\section{Searching Sorted Data}

Binary search is also used to search sorted lists.

Suppose we have alphabetically sorted names:

\begin{center}
Alice \\
Charlie \\
Dave \\
Frieda \\
Kallie \\
Moe \\
Paul \\
Sarah \\
Xavier
\end{center}

We repeatedly check the middle item and eliminate half of the remaining list.

\section{Algorithm Efficiency}

Algorithm efficiency describes how runtime scales with input size.

This is usually expressed using Big-O notation.

Common examples:

\[
O(n)
\]

\[
O(n^2)
\]

\[
O(\log n)
\]

\section{Linear Search}

Searching a list one item at a time requires approximately:

\[
\frac{n}{2}
\]

checks on average.

Complexity:

\[
O(n)
\]

\section{Binary Search Efficiency}

Binary search repeatedly cuts the search space in half.

If:

\[
n = 2^h
\]

then:

\[
h = \log_2(n)
\]

Therefore binary search complexity is:

\[
O(\log n)
\]

\section{Why Logarithmic Complexity Is Powerful}

Suppose:

\[
n = 100,000,000
\]

A linear algorithm may require roughly:

\[
100,000,000
\]

operations.

But binary search only requires approximately:

\[
\log_2(100,000,000) \approx 27
\]

steps.

This is dramatically faster.

\section{Key Takeaways}

\begin{itemize}
    \item Recursive functions call themselves.
    \item Recursive algorithms require stopping conditions.
    \item Binary search repeatedly divides a problem in half.
    \item Binary search is extremely efficient.
    \item Big-O notation describes scaling behavior.
    \item Logarithmic algorithms are vastly faster than linear algorithms for large datasets.
\end{itemize}

\end{document}
